\section{Conclusion and Outlook}\label{sec:conclusion}

We have presented a framework to solve the central dilemma of long-document processing: how to scale semantic extraction to massive corpora without losing the ability to verify the results. By treating hierarchical summarization not as a black-box compression but as an algebraic reduction, we allow practitioners to audit the fidelity of the entire pipeline by inspecting only its atomic parts.

Our approach rests on a semantic generalization of \emph{mergeable summaries}. While the database literature relies on explicit, commutative sketches to preserve counts and frequencies, we substitute these with ``learned'' consistency conditions---Sufficiency \ref{eq:leaf_sufficiency}, Idempotence \ref{eq:leaf_idempotence}, and Merge Consistency \ref{eq:leaf_compatibility}---that force stochastic LLMs to behave \emph{as if} they were rigorous algebraic operators. The LLM never becomes truly associative; it merely mimics the merge algebra closely enough that the oracle cannot distinguish the two routes. This connects the flexibility of modern language models with the structural guarantees of randomized streaming algorithms.

The primary contribution is operational. We replace the impossible requirement of end-to-end verification (reading the whole library to check the summary) with a scalable \emph{probabilistic audit}. By sampling realized leaves and internal merges, estimating their failure rates, and monitoring especially the leaf violations that compare summaries to raw text, practitioners obtain a live picture of how often the hierarchy corrupts the oracle. These empirical error rates make visible the precise trade-offs between block size, merge schedule, and accuracy, allowing social scientists to certify the reliability of their data extraction without needing to manually code a statistically significant fraction of the full corpus. The guarantees are about process quality: if the measured per-edge failure rates are small, the pipeline is trustworthy; for any specific document, reliability still follows the familiar $(1-p)^N$ scaling.

These guarantees extend beyond measurement to learning. We showed that under standard factorization assumptions (where supervision depends only on the oracle $\fstar$), training Direct Preference Optimization (DPO) policies on validated summaries is information-equivalent to training on full documents. The gap between the ideal policy and the learned policy is controlled by the empirically estimated violation rates, so continued sampling tightens the alignment guarantee. Because the oracle can encode arbitrarily nuanced researcher preferences (provided they are expensive but finite to evaluate), this route gives a practical way to align LLM-generated summaries with whatever semantics matter, while explicitly acknowledging that stylistic alignment outside the oracle's scope must be handled separately.

The limits of this approach are clear. The strongest equivalences hold only when supervision strictly factors through the task oracle; if human annotators reward presentational features lost during summarization, the alignment may drift. Likewise, substitution consistency (Condition~\eqref{eq:leaf_compatibility}) can fail if summaries drop context that later merges rely on. Those cases show up as merge-consistency violations in the audit and must be addressed by prompt tuning or stronger models. Finally, on-range stability is substantive: without it, additional ``clean-up'' rounds can silently flip labels, a failure mode our audit is designed to detect.

Methodologically, the blueprint is simple but demanding: Build hierarchies that pass the edge-wise audit; optimize models to satisfy merge consistency; and report the estimated violation rates alongside codebooks and regression tables. The result is a pipeline whose guarantees align with the research design, whose failures are interpretable, and whose outputs can be trusted to mean what the social scientist claims they mean.

\bibliographystyle{plainnat}
\bibliography{refs}
